\capituloPregunta{¿Cómo construir un proyecto experimental completo?}

\section{Motivación}

Un proyecto experimental no es una suma de cálculos. Es una cadena de decisiones: problema, diseño, medición, análisis, interpretación y comunicación.

\section{Materiales integrados}

La base del curso contiene proyectos sobre acústica, compósitos, inductancia, rendimiento térmico y productos alimentarios. Estos materiales deben usarse como modelos para mostrar cómo una pregunta técnica se transforma en un reporte defendible.

\section{Estructura mínima del proyecto}

\begin{longtable}{p{0.28\textwidth}p{0.62\textwidth}}
\toprule
Sección & Pregunta que debe responder\\
\midrule
Planteamiento & ¿Qué problema de ingeniería se quiere resolver?\\
Objetivo & ¿Qué decisión se busca tomar?\\
Diseño & ¿Qué factores, niveles y respuestas se estudiarán?\\
Datos & ¿Cómo se obtuvieron y qué limitaciones tienen?\\
Análisis & ¿Qué método corresponde al diseño?\\
Resultados & ¿Qué evidencia se obtuvo?\\
Discusión & ¿Qué se puede y qué no se puede concluir?\\
Conclusión & ¿Qué decisión se recomienda?\\
\bottomrule
\end{longtable}

\section{Actividad final}

Elija una práctica del manual y conviértala en propuesta de proyecto. Debe incluir:

\begin{enumerate}
  \item pregunta experimental;
  \item justificación de ingeniería;
  \item factores y niveles;
  \item unidad experimental;
  \item variables respuesta;
  \item plan de análisis;
  \item riesgos y limitaciones;
  \item forma de comunicar resultados.
\end{enumerate}

\begin{cuidado}
Un proyecto no debe concluir más de lo que su diseño permite. La honestidad sobre límites es parte del rigor.
\end{cuidado}
